The signature of a continuous nondegenerate symmetric form on a connected base is constant: the eigenvalues of its Gram matrix with respect to any auxiliary frame metric vary continuously and never vanish, so their sign counts are locally constant, and $M$ is connected. Suppose $h$ nondegenerate everywhere. On the physical exterior $h = G \perp K$ has signature $(p,q) + \operatorname{sig}(K)$; on the far exterior it has $(p{-}1, q{+}1) + \operatorname{sig}(K)$ (Remark \ref{rem:coords}, the wall being tame, \claimref{sheaf.tameness}), with the \emph{same} $\operatorname{sig}(K)$ by the standardness hypothesis. The two totals differ in the first entry by one; a single form of constant signature cannot have both, so nondegeneracy must fail somewhere in $U$. The constancy clause on $\operatorname{sig}(K)$ is necessary: $M \setminus U$ is disconnected (the two sides lie in different components of $M \setminus \Sig$), and a completion allowed to switch $K$'s signature across the wall, say from $(1,1)$ to $(2,0)$, could match the two totals; ``standard'' means one exterior model, not two.
